\subsection*{如何看待并行任务}

在接触 Croqtile 语法之前，先把「并行」到底是什么意思说清楚。

假设有八个彼此独立的任务——八个要相加的分块、八行要处理的数据，等。你可以一次把八个任务都描述出来，称之为「并行」。但这句话完全没说明有多少任务真的在同时运行：在单核 CPU 上，它们一个接一个执行；在 4 核 CPU 上，也许同时执行四个。

\begin{figure}[htbp]\centering
\includegraphics[width=\linewidth,height=0.56\textheight,keepaspectratio]{f05-fig1_virtual_parallelism_light.png}
\caption{同样的八个任务，在每次执行一个、四个或八个任务的示意调度下，分别需要八个、两个或一个时间步。}
\end{figure}

并行性描述哪些任务可以独立执行，执行资源则决定同一时刻能推进多少任务。

如果分块数量和启动尺寸由同一组常量表示，修改分块就会迫使资源映射一起改变。将两种描述分开后，这些选择更容易调整。

Croqtile 用 \co@parallel@ 描述独立工作，用空间说明符描述映射。因此，可以先理解工作划分，再判断某个轴应映射到块、线程还是协作组。
